% ====================================================================
\section{Correspondence with the original manuscript}\label{app:map}
% ====================================================================

The table maps every section of the original manuscript~\cite{Src} to its place in this version.
Nothing of mathematical substance was removed. The appendix on source files and internal review status,
which is not mathematical, is summarised at the end of this appendix.

\begin{center}\small
\begin{longtable}{p{4.6cm}p{4.6cm}p{5.2cm}}\toprule
original & this version & change\\\midrule\endhead
\S1 model, Def.~1.1, Thm.~1.2 & \cref{sec:prelim}, \cref{def:problems}, \cref{thmB,thmC} & split into model and results\\
\S1 source problem (ordered MCVP) & \cref{sec:model,sec:pcomplete} & recast as MSLPs; classical references made explicit; normal-form lemma added\\
\S2.1 robust generators, constant sources & \cref{def:robust}, \cref{sec:timers} & feedback made explicit in the definition\\
Lemma 2.1 (timers) & \cref{lem:timers} & unchanged\\
Lemma 2.2 (pulses) & \cref{lem:pulses} & unchanged\\
Thm.~2.3 (phase defects) & \cref{thm:phasedefect} & unchanged\\
\S2.4 combiner & \cref{lem:combiner} & three cases merged into two\\
\S2.5 bases, size & proof of \cref{thm:phasedefect}, \cref{lem:paritytransform} & parity transformation isolated as a lemma\\
\S2.6 logspace enumeration & \cref{lem:genlogspace} & unchanged\\
\S2.7 realised words & \cref{lem:charwords}, \cref{cor:words} & unchanged\\
\S3.1 \textsf{EDGE} cell, Lemma 3.1 & \cref{lem:modes,lem:forcing,lem:copy} & forcing hypothesis weakened (\cref{rem:forcing-hyp})\\
\S3.2 move & \cref{lem:move}, \cref{rem:nopause} & unchanged\\
\S3.3 junction, routing & \cref{lem:route}, \cref{lem:frame} & frame lemma factored out; chain transport \cref{lem:chain} added\\
\S3.4 gates & \cref{lem:gate} & unchanged\\
\S3.5 phase assignment & \cref{sec:phasecalc}, \cref{def:procphases} & cell-phase, concatenation and timing lemmas\\
\S4 read (Lemma 4.1), \S4.2 & \cref{lem:read}, \cref{tab:timing} & unchanged; timing via the calculus\\
\S5.1 processor, normal form & \cref{def:processor,def:normal,lem:proctopology,lem:normalstable} & stability made a lemma\\
\S5.2 alignment waits & \cref{lem:sigmatiming,lem:waits} & waits enumerated; stability of each waiting configuration proved\\
\S5.3 prepare, commit & \cref{lem:prepare,lem:commit} & unchanged\\
\S5.4 fetch, store, recover & \cref{lem:fetch,lem:store,lem:recover} & fetch proved through \cref{lem:chain}\\
\S5.5 gate, cleanup, routes & \cref{lem:gatestep,lem:clean,lem:writeroutes}, \cref{rem:fronts} & unchanged\\
\S5.6 normalisation & \cref{lem:normalize} & wave invariant $W_j$ made explicit\\
\S5.7 program invariant & \cref{prop:instruction,thm:forced}, \cref{tab:milestones} & milestones tabulated; closed formula for $\tau$ (\cref{cor:tau})\\
\S6.1 control replacement & \cref{prop:substitution,prop:correct} & substitution principle stated and proved\\
\S6.2 size and time & \cref{prop:size}, \cref{cor:tau} & horizon $\Hz=\tau+3$, or any larger polynomial\\
\S6.3 logspace transducer & \cref{prop:logspace} & unchanged\\
\S7 membership in $\Ptime$ & \cref{prop:membership} & attributed to the classical energy method\\
\S8.1 reader (Lemma 8.1) & \cref{lem:reader} & unchanged\\
\S8.2 final write & \cref{con:ever}, \cref{lem:finalwrite} & unchanged\\
Scope paragraph & \cref{rem:scope} & unchanged\\\bottomrule
\end{longtable}
\end{center}

\paragraph{Provenance.} The original English edition assembles research notes dated 8 October
2026. Those notes record internal readings of the reduction and are attributed there to automated
research agents. No human authorship, publication status, independent external review or
proof-assistant verification is asserted there, and none is asserted here. The present version is a
reorganisation; the computational checks of \cref{app:audit} are sanity checks, not part of the
proof.

% ====================================================================
\section{Audit notes}\label{app:audit}
% ====================================================================

\paragraph{Points made explicit by the reorganisation.}
\begin{enumerate}[label=(\arabic*)]
\item \emph{Duality.} ``Complementing the contract'' is formalised at the level of modes
  (\cref{lem:duality}). Literal complementation of control bits would replace $\HOLD=(0,1)$ by
  $(1,0)$ (\cref{rem:dualcontrols}). Duality shifts the admissible start parity by one, which is the
  background term $c$ in \eqref{eq:startparity}.
\item \emph{Merging of control blocks.} The phase rule of the original requires each mode block to
  start at the right parity and to have odd length. It does not address two adjacent blocks of the
  same mode, which would merge into one block of even length and violate \eqref{eq:phase}.
  \Cref{lem:concat} proves that this cannot happen when all blocks are admissible.
\item \emph{Alignment waits.} The original permits ``a zero- or one-transition wait whenever the
  preparation is stable''. Here the waits are determined exactly (\cref{lem:sigmatiming}): six per
  instruction plus one after each normalisation. The stability of every waiting configuration is
  proved (\cref{lem:waits}), and the observation time follows in closed form (\cref{cor:tau}).
\item \emph{Feedback through consumers.} Robustness quantifies over all consumer histories, so it
  covers feedback. The substitution principle (\cref{prop:substitution}) makes the induction
  explicit.
\item \emph{Frame conditions.} Every junction outside the footprint of a primitive is checked
  through \cref{lem:frame} in the corresponding proof. For bus transports, the far ports of
  intermediate destinations are handled once and for all inside \cref{lem:chain}.
\item \emph{Space bounds.} The space of the reduction is kept separate from the complexity of
  evaluation (\cref{rem:twospaces}).
\end{enumerate}
No error affecting the validity of the original argument was found: every schedule of the original
is reproduced unchanged.

\paragraph{Computational sanity checks.} Two scripts accompany this version, in
\texttt{tools/check/}. They re-implement the constructions from the text.
\begin{itemize}
\item \texttt{processor\_sim.py} simulates the processor at vertex level, with every control slot
  forced to its scheduled word. It checks:
  \begin{itemize}
  \item every milestone $\NORMAL(c_k,\rho_k)$ and the output;
  \item condition \eqref{eq:phase} for every control word, with its assigned phase;
  \item the absence of overlapping mode assignments;
  \item the formula $D_k=12(i+j)+6m+92$;
  \item the activation variant, including arbitrary perturbations of $R_m$ after the window.
  \end{itemize}
  It runs on all programs in normal form with at most three source nodes, plus the appended final
  gate, and on random programs with up to ten registers. Deliberate mutations are detected: omitting
  $\prim{SET}(F_A\mathbin{:=}c)$ in normalisation, a pause inside a move, removing the internal
  wait of a fetch, a read that does not restore $C$, and restoring the fronts before $E_0=c$.
\item \texttt{clock\_sim.py} builds the timers, the pulses and the phase-defect generators exactly as
  described, and simulates them against adaptive and adversarial consumers. It checks the degrees,
  that each output is a tree, the timer sizes, the counting recurrence \eqref{eq:counting} at every
  skeleton node, \cref{lem:charwords}, and \cref{lem:reader}. The generators are checked for all defect
  sets with dates at most $11$, in both phases, and for random larger ones.
\end{itemize}
These checks support the written proofs but do not replace them. The complete trees, of size
$O(m^{13})$, are not simulated end to end. The link between the forced processor and the
generators is \cref{prop:substitution}, which is proved, not tested.
